%    Latex Template for ACA 2015
%%%%%%%%%%%%%%%%%%%%%%%%%%%%%%%%%%%%%%%%
%    DO NOT CHANGE THE SETTINGS BELOW
%%%%%%%%%%%%%%%%%%%%%%%%%%%%%%%%%%%%%%%%
\documentclass[11pt]{article}
\usepackage{epsfig}
\usepackage{graphicx}
\usepackage[latin1]{inputenc}
\usepackage[T1]{fontenc}
\usepackage{mathptmx}
\usepackage{url}



\begin{document}



%%%%%%%%%%%%%%%%%%%%%%%%
%%%%%%%%%%%%%%%%%%%%%%%%%%%%%%%%%%%%%%%%
%    DO NOT CHANGE THE SETTINGS ABOVE
%%%%%%%%%%%%%%%%%%%%%%%%%%%%%%%%%%%%%%%%
%
\vspace{-0.8cm}
\begin{flushleft}
\Large \textbf{\noindent
%%% Supporting Mathematical Thinking with CAS: The Need of Epistemic Change among Teachers.
 %%%%%%%
Supporting Mathematical Thinking with CAS: The Need of Epistemic Change among Teachers.
}
%%%%%%%%%%%%%%%%%%%%%%%%%%%%%%%%%%%%%%%%%%%%%%%%%%%%%%%%
\\
%%%%%%%%%%%%%%%%%%%%%%%%%%%%%%
\vspace{0.5cm}
\normalsize
Rotem Abdu
 %%%%%% NAME OF SPEAKER SHOULD BE UNDERLINED %%%%%%%%%
\normalsize{
 %%%%%% FOR EXAMPLE %%%%%
% A. Maths$^1$, B. Sciences$^2$, \underline{C. Lemma}$^3$, D. Technology$^4$
 %%%%%%%%%%%%%%%%%%%%%%%%
} \\
\vspace{5mm}
\textit{\footnotesize
Levinsky College, Tel-Aviv, Israel
%%%%%%%%%%%%%%%%%%%%%%%%%%%%%%%%%%%%%
}
\end{flushleft}

     The current abundance of educational technologies, and computer algebra systems (CAS) in particular, carry a promise: new venues for advanced mathematical thinking. This promise is a product of teachers' ability to construct and simulate mathematical ideas that are dynamic and constrained by the mathematical world (represented with the CAS). This dynamic attribute is considered to bring a change in the way mathematical ideas are thought of: instead of prototypical examples that are drawn with pen and paper, CAS such as Geogebra can provide a dynamic context for inquiry of a full range of examples for mathematical concepts. Ultimately the goal is to create cohesive mental models of mathematical ideas. 

     The promise for new mathematical thinking afforded by CAS, however, is impeded by several factors. For example, Anthony and Clark, (2011) examine key factors that cause teachers to refrain from using CAS in their classrooms, including dilemmas of misalignment with other curricular goals and limited professional development. My interest in this talk, however, is in the change of epistemic stance that is required by teachers; from static prototypes to dynamic and invariant objects. 

     I will examine this point with a case study in which nine in-service teachers participated in an activity, in which they were asked to construct invariant models of mathematical objects --- right triangle --- with Geogebra. This activity was a part of a course on methodological issues in mathematics education. The teachers were familiar with Geogebra as part of their postgraduate curriculum. An analysis of the activity shows that the mathematical objects that were constructed by the teacher, at first, neglected the much necessary ``invariance'' attribute of such object: in the case of the right triangle created in Geogebra should stay a right triangle even if one of its vertices or segments is moved by the user. Upon instruction and refinement of the objective of the activity, teachers were gradually able to construct objects that are invariant. However, the idea of invariance became what Cobb et al., (2001) would call ``a socio-mathematical norm'' among the teachers, only after two more similar activities. 

I conclude that there is a need for epistemic change --- even among teachers that are supposedly familiar with Geogebra --- from seeing CAS as tools that afford static prototypes to seeing them as environments for building and simulating invariant-dynamic objects. For that matter, teachers need to participate in activities that would provide them with opportunities to make that epistemic change; namely, they should be engaged in building such models as well as observing other who do so. These results has also carry a nesting effect: if a teacher is not familiar with the invariant principle, there are good chances that their students will not adopt this epistemic stance either. Moreover, the illustrated case also suggest that ``frontal'' teaching with CAS --- without students' or teachers' hands-on experience and building of objects --- will yield limited learning in terms of achieving advanced mathematical thinking.

%\begin{eqnarray}
%\lim_{x \to +\infty} {\rm Abstract}(x) = 1\ {\rm page}.
%\label{eq1}
%\end{eqnarray}
%Your abstract should be prepared in \LaTeX\ according to the template provided  and then \textbf{converted to PDF for submission.}   Abstracts not complying with this template may not be included in the book of abstracts of the ACA 2015. Abstract submission should be done  via the session organizers, they will confirm your submission. If your abstract is accepted, you will be invited to present your talk.
%The book of abstracts will be available to the participants prior to the conference and will provide the audience with more information
%about the papers being presented at the ACA 2015 Workshop. At least one author must be registered for your talk to be included in the program of the conference.

%%%%%%%%%%%%%%%%%%%%%%%%%%%%%%%%%%%%%%%


%%%%%%%%%%%%%%%%%%%%%%%%%%%%%%%%%%%%%%%%
%% BIBLIOGRAPHY
%%%%%%%%%%%%%%%%%%%%%%%%%%%%%%%%%%%%%%%%
\begin{thebibliography}{00}
\addcontentsline{toc}{chapter}{References}
\setlength{\itemsep}{-1mm}
\footnotesize
\bibitem{Anthony} Anthony, A. B. and Clark, L. M., \textit{Examining dilemmas of practice associated with the integration of technology into mathematics classrooms serving urban students},�\textit{Urban Education},�\textbf{46}(6), 1300--1331 (2011).
\bibitem{Cobb} Cobb, P., Stephan, M., McClain, K., and Gravemeijer, K., \textit{Participating in classroom mathematical practices}, In�\textit{A journey in mathematics education research}�(pp. 117--163), Springer Netherlands (2010).

%\bibitem{KR} B. Keller and I. Reiten, \textit{Cluster-titled algebras are Gorenstein and stably Calabi-Yau}, Adv. Math. \textbf{211}, 1, pp. 123-151 (2007).
%\bibitem{WK} S. Washito and K. Karabashi, \textit{Use This Template}, in \textit{Proceedings of the 5th
%International Conference on Photon Interaction} (ICPI-2012), San Francisco, USA, ed. A. Chen, pp. 22-24 (2012).
%\bibitem{ABS} A.B. Smith, \textit{Describing the style for the AMMCS-CAIMS-2015 book of abstracts}, in C.D. Science (ed.), General Book of Wisdom, 2\textsuperscript{nd} ed., University of Queensland, pp. 23-42 (2011).
\end{thebibliography}
\normalsize
\end{document} 
